\begin{EnChapter}

\chapter{System Model}
\label{chap:system-model}

This chapter models downlink resource control in a multi-hop integrated access
and backhaul (IAB) network. The model permits arbitrary finite tree depth, unequal
branch depths, and different numbers of IAB nodes and user equipments (UEs).
Each IAB distributed unit (DU) retains its slot-level proportional fair (PF)
scheduler. A local deep reinforcement learning (DRL) agent adjusts the scheduling
eligibility of its directly attached UEs through a joint time-domain mask action;
the mobile terminations (MTs) of downstream IAB nodes remain eligible. Local
learning uses a deep Q-network (DQN) with a subtree-throughput reward. The
federated configuration shares model parameters through sample-weighted federated
averaging (FedAvg). Only PF, pure Local DRL, and Local DRL with FedAvg are defined
here. Numerical testbed settings are separated from the general formulation.

\section{Network Topology}
\label{subsec:topology}

Let $\Nset=\{n_1,\ldots,n_N\}$ be the IAB nodes and
$\Uset=\{u_1,\ldots,u_U\}$ the terminal UEs. The donor is denoted by $0$, and
these three sets of identities are disjoint. Define
\begin{equation}
  \mathbb{T}=(\Vset,\Eset),\qquad
  \Vset=\{0\}\cup\Nset\cup\Uset,\qquad
  \Eset=\{(\parent(v),v):v\in\Nset\cup\Uset\},
\end{equation}
where $\parent(v)\in\{0\}\cup\Nset$ is the unique serving parent of $v$.
The graph is a rooted directed tree with donor $0$ as its root and UEs as leaves.
Thus every UE has a unique downlink route. Multi-parent routing and topology
reconfiguration during a measurement run are outside the present model.

\begin{definition}[Children, depth, and descendants]
For $n\in\{0\}\cup\Nset$, define
\begin{align}
  \Cset_n^{\mathrm{I}}&=\{m\in\Nset:\parent(m)=n\}, &
  \Cset_n^{\mathrm{U}}&=\{u\in\Uset:\parent(u)=n\},\\
  \Cset_n&=\Cset_n^{\mathrm{I}}\cup\Cset_n^{\mathrm{U}}, &
  h(v)&=h(\parent(v))+1,\quad h(0)=0.
\end{align}
The maximum UE hop count is $H=\max_{u\in\Uset}h(u)$. Let $\Uset_v$ be the
terminal UEs descended from vertex $v$, including $v$ when $v$ is itself a UE.
Then
\begin{equation}
  \Uset_u=\{u\},\qquad
  \Uset_n=\Cset_n^{\mathrm{U}}\cup
           \bigcup_{m\in\Cset_n^{\mathrm{I}}}\Uset_m.
  \label{eq:subtree-ues}
\end{equation}
The union in~\eqref{eq:subtree-ues} is disjoint. The ordered downlink path
$\mathcal{P}(u)$ is the sequence of edges from $0$ to $u$.
\end{definition}

Any IAB node may serve both terminal UEs and downstream IAB nodes. A node with
$\Cset_n^{\mathrm{I}}\ne\emptyset$ is called a relay; an IAB node without IAB
children is called an access node. These roles depend on connectivity rather than
on a fixed depth or node identifier. There is no fixed bound on $N$, $U$, $H$, or
$|\Cset_n|$ in the mathematical model.

\paragraph{MT--DU architecture and forwarding.}
Node $n\in\Nset$ contains MT$_n$, attached to DU$_{\parent(n)}$, and DU$_n$,
serving $\Cset_n$. An IAB child appears to its parent's scheduler as one MT
scheduling entity, although that entity carries traffic for every UE in its
subtree. The donor connects to the core network through a wired interface.
The implementation forwards traffic using Linux IP routing over MT PDU sessions;
it does not implement a Backhaul Adaptation Protocol layer. Each child has one
configured data bearer, and the parent observes an aggregated RLC queue towards
that child. The per-flow queues below are an analytical decomposition of these
aggregates, rather than separately exposed queues at the controller.

\begin{example}[Current experimental instance]
The deployed topology has $N=12$, $U=24$, and $H=3$. The donor serves relays
$n_1,\ldots,n_4$. For $i=1,\ldots,4$,
\begin{equation}
  \Cset_{n_i}^{\mathrm{I}}=\{n_{2i+3},n_{2i+4}\},\qquad
  \Cset_{n_i}^{\mathrm{U}}=\{u_{2i+15},u_{2i+16}\}.
\end{equation}
For $i=5,\ldots,12$,
\begin{equation}
  \Cset_{n_i}^{\mathrm{I}}=\emptyset,\qquad
  \Cset_{n_i}^{\mathrm{U}}=\{u_{2i-9},u_{2i-8}\}.
\end{equation}
UEs $u_1,\ldots,u_{16}$ are three hops from the donor; relay-attached UEs
$u_{17},\ldots,u_{24}$ are two hops away. These counts instantiate the model and
are not constants in its queue, reward, or aggregation equations.
\end{example}

\section{Time Structure and Radio Resources}
\label{subsec:resource}

A radio slot is indexed by $t\in\mathbb{Z}_{\ge0}$ and has duration
$T_{\mathrm{slot}}=10^{-3}2^{-\mu}$ seconds, where $\mu$ is the numerology.
All times and rates in this chapter use simulated time unless explicitly stated
otherwise. DU$_n$ has $B_n$ PRBs. Let $\tilde B_n(t)\le B_n$ be the available
PRBs in slot $t$, with $\tilde B_n(t)=0$ when no downlink data can be scheduled.
The set $\Tset_n^{\mathrm{DL}}$ includes downlink-capable slots, including usable
special slots of the configured TDD pattern.

An observation is collected every $\Delta_o$ seconds. A normal decision interval
contains $L_h$ observations, with duration $\Delta_d=L_h\Delta_o$. Decision $k$
begins at slot $t_k$ and is held over $\Tset_k=\{t_k,\ldots,t_{k+1}-1\}$.
A change in the attached-child set may end the hold early and trigger a new
valid action. Node decisions need not be perfectly synchronized; timestamped
observations are aligned during training.

\begin{assumption}[Inter-link radio independence]
\label{ass:orthogonal}
The model assumes out-of-band-like operation: a node's MT reception does not
exclude simultaneous DU transmission, and different DU radio links have no
modelled co-channel interference. On the OAI RF-simulation platform the links use
independent connections with no cross-link interference, even though their
nominal carrier configurations can coincide. This behaviour is not evidence that
all deployed cells use physically non-overlapping carriers. No in-band MT--DU
half-duplex constraint is imposed in this study.
\end{assumption}

The earlier MT-activity-based DU PRB budget reduction is disabled. Consequently,
in a downlink-capable slot the current resource configuration has
\begin{equation}
  \tilde B_n(t)=\lfloor\beta_n(t)B_n\rfloor,\qquad
  \beta_n(t)=1,
  \label{eq:beta}
\end{equation}
subject to the scheduler's ordinary control, reference-signal, and allocation
constraints. Resource competition still exists among the children served by the
\emph{same} DU, and across hops through their traffic and queues.

\section{Link Service and PF Scheduling}
\label{subsec:rate}
\label{subsec:coupling}

Let $q_{n,v}(t)$ be the MCS selected for child $v\in\Cset_n$.
An approximate payload capacity in bits per PRB is
\begin{equation}
  \rho_{n,v}(t)=\nu_{n,v}(t)N_{\mathrm{RE},n,v}^{\mathrm{data}}(t)
               Q_m\bigl(q_{n,v}(t)\bigr)R\bigl(q_{n,v}(t)\bigr),
  \label{eq:rho}
\end{equation}
where $\nu$ is the layer count, $N_{\mathrm{RE}}^{\mathrm{data}}$ is the number of
payload resource elements per PRB, $Q_m$ is the modulation order, and $R$ is the
coding rate. Actual transport-block sizing is performed by OAI; this expression
is a service abstraction, not an application-throughput prediction. Let
$\eta_{n,v}(t)\in[0,1]$ represent the effective successfully delivered payload
fraction, including the effects of retransmissions and protocol overhead.

Let $b_{n,v}(t)$ be the total allocated PRBs, including new transmissions and
HARQ retransmissions. Let $b_{n,v}^{\mathrm{new}}(t)\le b_{n,v}(t)$ be the
new-transmission allocation and $z_{n,v}(t)=\Ind[b_{n,v}^{\mathrm{new}}(t)>0]$.
Ordinary DU constraints include
\begin{equation}
  b_{n,v}(t)\in\mathbb{Z}_{\ge0},\qquad
  \sum_{v\in\Cset_n}b_{n,v}(t)\le\tilde B_n(t),\qquad
  z_{n,v}(t)\le e_{n,v}(t),
  \label{eq:du-budget}
\end{equation}
where $e_{n,v}(t)$ is the new-transmission eligibility indicator, accounting for
attachment, HARQ processes, feedback opportunities, and other MAC constraints.
For example, a finite feedback opportunity $g$ can constrain scheduled
transmissions through
\begin{equation}
  \sum_{t\in\mathcal{F}_{n,v,g}}z_{n,v}(t)\le A_{n,v,g},
  \label{eq:feedback}
\end{equation}
where $\mathcal{F}_{n,v,g}$ contains slots mapped to that opportunity and
$A_{n,v,g}$ is its feedback capacity. This is an abstraction of the configured
feedback mechanism, not a universal fixed transmission limit for all base stations.
Unused PRBs can remain when no other child is eligible; PRB availability alone
therefore does not guarantee that another child can receive them.

PF orders eligible children using their estimated instantaneous supportable rate
relative to their smoothed served rate, abstractly
\begin{equation}
  M_{n,v}(t)=\frac{\widehat R_{n,v}^{\mathrm{inst}}(t)}
                    {\overline R_{n,v}(t)+\epsilon_{\mathrm{PF}}}.
  \label{eq:pf-metric}
\end{equation}
Retransmission allocations follow the existing HARQ rules and consume the same
DU budget. The existing OAI scheduler performs the resulting PRB allocations. DRL does not
replace this scheduler, choose MCS values, or directly assign individual PRBs.

\section{Joint Time-Domain Control}
\label{subsec:action}

Let $\mathcal{M}$ be a finite set of binary masks of width $W$, containing the
all-open mask $m_{\mathrm{PF}}$. Let $\Cset_n^{\mathrm{att}}(k)$ be the children
reported as attached at decision time, including attached children with empty
queues. The eligible mask targets are
\begin{equation}
  \mathcal{D}_n(k)=\Cset_n^{\mathrm{att}}(k)\cap\Cset_n^{\mathrm{U}}.
\end{equation}
The node chooses a \emph{joint} action: a subset of direct UEs and one common mask
for that subset. The candidate set is
\begin{equation}
  \Aset_n(o_n(k))=\{a_{\mathrm{PF}}\}\cup
  \bigl\{(\mathcal{S},m):\emptyset\ne\mathcal{S}\subseteq\mathcal{D}_n(k),
                         m\in\mathcal{M}\setminus\{m_{\mathrm{PF}}\}\bigr\},
  \label{eq:action-set}
\end{equation}
with cardinality
\begin{equation}
  |\Aset_n|=1+(2^{|\mathcal{D}_n|}-1)(|\mathcal{M}|-1).
  \label{eq:action-count}
\end{equation}
The PF action is represented once, independently of subset membership.
For $a_n(k)=(\mathcal{S}_n(k),m_n(k))$, define
\begin{equation}
  m_{n,v}(k)=
  \begin{cases}
    m_n(k),&v\in\mathcal{S}_n(k),\\
    m_{\mathrm{PF}},&v\notin\mathcal{S}_n(k).
  \end{cases}
  \label{eq:child-mask}
\end{equation}
For $a_{\mathrm{PF}}$, every child receives $m_{\mathrm{PF}}$.
In particular, $m_{n,m}=m_{\mathrm{PF}}$ for all IAB children
$m\in\Cset_n^{\mathrm{I}}$: the current policy never masks a downstream MT.

Let $\iota_n(t)\in\{0,\ldots,W-1\}$ be the mask-bit index used by the MAC.
The additional scheduling constraint is
\begin{equation}
  z_{n,v}(t)\le m_{n,v}(k)[\iota_n(t)],\qquad t\in\Tset_k,
  \label{eq:mask-constraint}
\end{equation}
where a one bit permits scheduling and a zero bit excludes new scheduling through
the controlled PF path. The implementation checks the mask in the new-transmission
scheduling path; HARQ retransmissions retain their own scheduling rules. Thus the
constraint models the controlled transmissions, rather than an assurance that no
radio transmission of any kind occurs in a masked slot.

The PRB cap is left at $B_n$ for every child, so PF determines frequency-domain
allocation within eligible slots. The donor also remains unmodified PF.
The mask is an eligibility restriction, not an exclusive slot reservation for
any other child. Its effective allowed downlink fraction is
\begin{equation}
  d_{n,v}(k)=
  \frac{\sum_{t\in\Tset_k}\Ind[t\in\Tset_n^{\mathrm{DL}}]
                          m_{n,v}(k)[\iota_n(t)]}
       {\max\{1,\sum_{t\in\Tset_k}\Ind[t\in\Tset_n^{\mathrm{DL}}]\}}.
  \label{eq:effective-mask}
\end{equation}
Masks with equal bit counts can have different effects because the TDD pattern,
special-slot symbol budgets, and feedback timing also matter. No universal
throughput gain is assumed for time-domain control relative to frequency control.

\section{Multi-Hop Traffic and Queues}
\label{subsec:traffic}

UE $u$ has an offered downlink rate $\lambda_u(k)\ge0$, in bits per simulated
second. Let $A_u(t)$ be new payload bits admitted at the donor. TCP injection is
adapted by transport congestion control; UDP can continue injecting its target
rate even when queues overflow. The offered rate is not an observation supplied
by the scenario to the DRL policy.

For every DU $n$ on $\mathcal{P}(u)$, let $Q_n^{(u)}(t)$ be the buffered bits of
flow $u$, and let $S_n^{(u)}(t)$ be successfully forwarded bits. For child $v$,
\begin{equation}
  Q_{n,v}(t)=\sum_{u\in\Uset_v}Q_n^{(u)}(t),\qquad
  S_{n,v}(t)=\sum_{u\in\Uset_v}S_n^{(u)}(t).
\end{equation}
An abstract aggregate service limit is
\begin{equation}
  0\le S_{n,v}(t)\le
  \min\{Q_{n,v}(t),\eta_{n,v}(t)b_{n,v}(t)\rho_{n,v}(t)\}.
  \label{eq:queue-agg}
\end{equation}
Per-hop arrival conservation is
\begin{equation}
  A_0^{(u)}(t)=A_u(t),\qquad
  A_n^{(u)}(t)=S_{\parent(n)}^{(u)}(t-\tau_n),\quad n\in\Nset,
  \label{eq:flow-conservation}
\end{equation}
where $\tau_n\ge1$ is an abstract forwarding delay in slots and service before
slot zero is zero. Arrivals are placed in the queue after the current-slot
service. With capacity $Q_{n,v}^{\max}$ and
$A_{n,v}=\sum_{u\in\Uset_v}A_n^{(u)}$, queue evolution and overflow are
\begin{align}
  Q_{n,v}(t+1)&=\min\{Q_{n,v}^{\max},\ Q_{n,v}(t)-S_{n,v}(t)+A_{n,v}(t)\},
  \label{eq:queue}\\
  D_{n,v}(t)&=[Q_{n,v}(t)-S_{n,v}(t)+A_{n,v}(t)-Q_{n,v}^{\max}]^+.
  \label{eq:drop}
\end{align}
These equations apply at every depth. Aggregate service is split among flows in
queue order; flow-level queues are bookkeeping variables for end-to-end accounting.
Finite buffers keep occupancies bounded, but bounded queues alone do not imply
that every offered flow is fully supported.

The radio-delivered terminal rate and its long-term constraints are
\begin{align}
  x_u^{\mathrm{rad}}(k)&=\frac{1}{\Delta_k}
                  \sum_{t\in\Tset_k}S_{\parent(u),u}(t),
                  \qquad \Delta_k=|\Tset_k|T_{\mathrm{slot}},
  \label{eq:ue-tput}\\
  \overline x_u^{\mathrm{rad}}&\le\overline\lambda_u, &
  \sum_{u\in\Uset_v}\overline x_u^{\mathrm{rad}}
      &\le\limsup_{T\to\infty}\frac{1}{TT_{\mathrm{slot}}}
                \sum_{t=0}^{T-1}\E[S_{\parent(v),v}(t)].
  \label{eq:link-cap}
\end{align}
These are necessary flow constraints, not a sufficient characterization of all
feasible allocations under HARQ, PF, and transport dynamics. Backhaul bytes are
intermediate traffic and must not be added again to terminal throughput.

\section{Optimization Objective and Evaluation}
\label{subsec:problem}

Let $x_u^{\mathrm{app}}(k)$ be the unique payload delivered to the UE application
per simulated second over interval $k$. It can differ from the radio rate because
of overhead, retransmissions, reordering, and measurement boundaries. The target
of the control policy is end-to-end sum goodput,
\begin{equation}
  \textbf{(P1)}\quad
  \max_{\pi}\ \frac{1}{\sum_{k=0}^{K-1}\Delta_k}
     \sum_{k=0}^{K-1}\Delta_k\,
         \E_\pi\!\left[\sum_{u\in\Uset}x_u^{\mathrm{app}}(k)\right],
  \label{eq:P1}
\end{equation}
subject to~\eqref{eq:du-budget},~\eqref{eq:action-set},
\eqref{eq:mask-constraint}, and~\eqref{eq:queue-agg}--\eqref{eq:drop}.
The donor uses $a_{\mathrm{PF}}$; the learned policy comprises only agents
$n\in\Nset$. For equal interval lengths this is the ordinary average sum goodput.
For congestion-only evaluation, the intervals in (P1) are restricted to the
predefined congestion measurement windows, with identical warm-up exclusions
for all methods. Scenario labels select evaluation windows and diagnostic
categories; they are not policy inputs.

Demand satisfaction and its fairness can be reported separately. For
$\Uset^+(k)=\{u:\lambda_u(k)>0\}$, define
\begin{equation}
  \sigma_u(k)=\min\{1,x_u^{\mathrm{app}}(k)/\lambda_u(k)\},\qquad
  J(k)=\frac{\bigl(\sum_{u\in\Uset^+(k)}\sigma_u(k)\bigr)^2}
             {|\Uset^+(k)|\sum_{u\in\Uset^+(k)}\sigma_u(k)^2}.
  \label{eq:satisfaction}
\end{equation}
$J(k)$ is undefined if there are no positive-demand UEs or all satisfaction ratios
are zero; such windows are identified rather than silently treated as perfect
fairness. The current Local DRL and FedAvg comparison optimizes throughput;
there is no implemented satisfaction constraint or Lagrangian fairness update
in this formulation.

Joint actions interact through downstream arrivals and transport feedback.
A higher served backhaul rate helps the objective only if the downstream path
ultimately delivers additional terminal payload. Restricting low-efficiency
direct UEs can increase subtree throughput when the downstream gain exceeds the
direct-UE loss; it can reduce throughput when downstream demand is already met.
Neither this gain nor a benefit from combining relay and access controls is
guaranteed by queue occupancy alone.

\section{Local DQN Formulation}
\label{subsec:marl}

\paragraph{Decentralized interaction.}
Each $n\in\Nset$ is an independently learning agent in a partially observable
multi-agent environment. The global state contains channels, queues, transport
state, and scheduler eligibility. The agent receives $o_n(k)$, selects
$a_n(k)\in\Aset_n(o_n(k))$, and later receives a subtree reward. Other agents'
policies can change during training, so the local environment is non-stationary.
Independent DQN and federated parameter sharing are approximations; they do not
provide a proof of optimality or convergence for (P1).

\paragraph{Observation.}
For a child $v$ reported as attached, let $Y_{n,v}^{\mathrm{prev}}(k)$ be the MAC downlink
transport-block bytes in the observation interval immediately preceding decision
$k$, $\widetilde Q_{n,v}(k)=Q_{n,v}(t_k)/8$ the reported queue in bytes, and
$q_{n,v}(k)$ the latest reported MCS. With reference values $Y_{\mathrm{ref}}$,
$Q_{\mathrm{ref}}$, and $q_{\mathrm{ref}}$, define
\begin{equation}
  \bm z_{n,v}(k)=\left[
  \frac{\log(1+Y_{n,v}^{\mathrm{prev}}(k))}{\log(1+Y_{\mathrm{ref}})},
  \frac{q_{n,v}(k)}{q_{\mathrm{ref}}},
  \frac{\log(1+\widetilde Q_{n,v}(k))}{\log(1+Q_{\mathrm{ref}})},
  \Ind[v\in\Cset_n^{\mathrm{I}}],
  \Ind[m_{n,v}(k-1)\ne m_{\mathrm{PF}}]\right].
  \label{eq:child-feature}
\end{equation}
The stored field named \texttt{wb\_cqi} carries MCS in this implementation, not
an independently measured CQI. Normalization references set the scale; log-scaled
features can exceed one if measurements exceed their references.

Let $L_j^{\mathrm{buf}}(k)$ be the sum of the latest RLC queue bytes reported by
DU$_j$. Neighbour observations give
\begin{align}
  \hat p_n(k)&=\clip\!\left(
       \frac{\log(1+L_{\parent(n)}^{\mathrm{buf}}(k))}
            {\log(1+Q_{p,\mathrm{ref}})},0,1\right),
  \label{eq:phat}\\
  \hat c_n(k)&=\clip\!\left(
       \frac{\log(1+\sum_{m\in\Cset_n^{\mathrm{I}}}L_m^{\mathrm{buf}}(k))}
            {\log(1+\max\{1,|\Cset_n^{\mathrm{I}}|\}Q_{\mathrm{ref}})},0,1\right).
  \label{eq:chat}
\end{align}
Only fresh available measurements enter these expressions. A missing or stale
parent measurement yields $\hat p_n=0$; no fresh IAB-child measurements yields
$\hat c_n=0$. Partial child information uses only the fresh children in the
numerator. $\hat p_n$ describes total parent congestion, not the specific parent
queue serving MT$_n$, so upstream congestion of another branch can also raise it.
For nodes directly below the donor, it is zero in the current implementation,
which does not collect donor relational data.

The complete observation is
\begin{equation}
  o_n(k)=\left(\{\bm z_{n,v}(k)\}_{v\in\Cset_n^{\mathrm{att}}(k)},\ \bm g_n(k)\right),
  \qquad
  \bm g_n(k)=\left[\frac{|\Cset_n^{\mathrm{att}}(k)|}{C_{\max}},
                  \widetilde\phi_n(k),\ \beta_n(k),\ \hat p_n(k),\ \hat c_n(k)\right].
  \label{eq:observation}
\end{equation}
The context $\widetilde\phi_n$ is the retained normalized Global xApp fairness
signal. When enabled, the signal compares recent mean node throughput with that
of nodes having the same topological role:
\begin{equation}
  \widetilde\phi_n=
  \frac{\clip(\overline T_{\mathrm{role}(n)}/(\overline T_n+\epsilon),0.5,2)-0.5}{1.5}.
  \label{eq:fairness-feature}
\end{equation}
When no signal is available, or the feature is disabled, its neutral value is
$1/3$. This is a soft observation feature, not a PRB budget or an FL weight.
Its enabled/neutral setting is identical in the Local-only and FedAvg experiments.
There is no proposed global bottleneck broadcast implemented by the equations here.

\paragraph{Subtree reward and measurement distinction.}
Let $\widehat x_u(\ell)$ be the terminal UE MAC throughput proxy in observation
interval $\ell$. Write $Y_{n,v}(\ell)$ for the MAC byte count collected
during that interval. The proxy is converted to Mbps per simulated second:
\begin{equation}
  \widehat x_u(\ell)=\frac{8\,Y_{\parent(u),u}(\ell)}{10^6\Delta_o}.
  \label{eq:mac-proxy}
\end{equation}
The general subtree utility and current reward are
\begin{align}
  r_n^{(\alpha)}(\ell)&=\sum_{u\in\Uset_n}U_\alpha(\widehat x_u(\ell)), &
  U_\alpha(x)&=
  \begin{cases}
    x^{1-\alpha}/(1-\alpha),&\alpha\ne1,\\
    \log(x+\epsilon_x),&\alpha=1,
  \end{cases}\\
  r_n(\ell)&=r_n^{(0)}(\ell)=\sum_{u\in\Uset_n}\widehat x_u(\ell).
  \label{eq:reward}
\end{align}
No MT throughput is added to this sum, preventing within-reward double counting
of forwarded terminal traffic. For arbitrary tree depth, the terminal set is
computed recursively by~\eqref{eq:subtree-ues}. In the current two-IAB-tier
implementation, a relay reward joins its own terminal-UE measurements with those
of its two access children. Timestamp joins outside the allowed tolerance are
discarded. Missing descendant measurements are not interpreted as zero reward.
Extending the implementation to deeper trees requires discovering all descendant
terminal-serving nodes and aligning their measurements.

The training reward is a MAC-derived proxy, whereas (P1) is assessed using
receiver-side application delivery. Also, rewards of ancestor and descendant
agents overlap: $\sum_{n\in\Nset}r_n$ is generally not the network sum throughput.
The use of subtree rewards encourages consideration of downstream delivery, but
neither makes local maximization exactly equivalent to (P1) nor removes
cross-branch effects.

\paragraph{Action holds and multi-step returns.}
The observation-level rewards belonging to one held decision are averaged:
\begin{equation}
  \overline r_n(k)=\frac{1}{L_{n,k}}
            \sum_{\ell\in\mathcal{H}_{n,k}}r_n(\ell),\qquad
            L_{n,k}=|\mathcal{H}_{n,k}|,
  \label{eq:hold-reward}
\end{equation}
where $\mathcal{H}_{n,k}$ contains the observations from that decision.
The merged transition retains the initial observation, the joint action, the
mean reward, and the observation after the hold. An incomplete group lacking
its initial decision record is discarded. Time-contiguous decision transitions
produce
\begin{equation}
  G_{n,k}^{(m)}=\sum_{j=0}^{m-1}\gamma^j\overline r_n(k+j),\qquad 1\le m\le n_{\mathrm{step}}.
  \label{eq:nstep-return}
\end{equation}
A discontinuity truncates the return. The discount $\gamma$ is per decision,
not per radio slot. The current pipeline stores a bootstrap discount $\gamma^m$
for a retained truncated sequence; truncation is not automatically an absorbing
terminal state. Phase-separated offline training and validation must construct
returns within their own partitions. A random split of overlapping online
multi-step transitions is only a training diagnostic, not an independent
performance test.

\paragraph{Q function and action selection.}
The online network $Q_{\theta_n}(o,a)$ estimates the discounted subtree return.
Its target network is $Q_{\theta_n^-}(o,a)$. The joint action is encoded by
per-child subset indicators and a mask-tier one-hot vector. In the current flat
network,
\begin{equation}
  Q_{\theta_n}(o,a)=c_Q f_{\theta_n}
       \bigl([\operatorname{pad}_{C_{\max}}(o),\ \bm g^{\mathrm{apply}}(a),
                   \operatorname{onehot}(m(a))]\bigr),
  \label{eq:q-network}
\end{equation}
where $c_Q>0$ rescales the network output and $\bm g^{\mathrm{apply}}$ is a
binary vector indicating the selected children. The flat MLP has two hidden
layers. It uses the same input layout at every node; it is not invariant to
permutations of child positions. A consistent child-to-position mapping is
therefore necessary, and permutation robustness must be measured rather than
assumed. Shared set or additive networks available as alternative code paths
are not substituted for the current flat network in this formulation.

The greedy candidate is
\begin{equation}
  a_n^*(k)=\operatorname*{arg\,max}_{a\in\Aset_n(o_n(k))}Q_{\theta_n}(o_n(k),a).
\end{equation}
An optional PF margin $\delta\ge0$, in return units, yields
\begin{equation}
  a_n^{\mathrm{greedy}}(k)=
  \begin{cases}
    a_n^*(k),&Q_{\theta_n}(o_n(k),a_n^*)-Q_{\theta_n}(o_n(k),a_{\mathrm{PF}})>\delta,\\
    a_{\mathrm{PF}},&\text{otherwise}.
  \end{cases}
  \label{eq:pf-margin}
\end{equation}
Training uses $\epsilon$-greedy exploration over valid joint actions, with a
configurable exploration distribution; evaluation sets $\epsilon=0$ and freezes
learning. A PF-only warm start is optional, while pretrained models can start
without it. The margin tests a learned return estimate, not a guaranteed goodput
improvement. The PF margin is an inference rule; the Double DQN target below
still maximizes over the full valid candidate set.

\paragraph{Double DQN update.}
For a replay transition with $m$ observed decisions, define
\begin{align}
  a'&=\operatorname*{arg\,max}_{a\in\Aset_n(o')}Q_{\theta_n}(o',a),\\
  y&=G_{n,k}^{(m)}+\gamma^m Q_{\theta_n^-}(o',a'),
  \label{eq:ddqn-target}\\
  \mathcal{L}_n(\theta_n)&=\frac{1}{|\mathcal{B}_n|}
          \sum_{(o,a,G,o',m)\in\mathcal{B}_n}
          \operatorname{Huber}\!\left(\frac{Q_{\theta_n}(o,a)-y}{c_Q}\right),
  \label{eq:dqn-loss}\\
  \theta_n&\leftarrow\operatorname{AdamStep}(\theta_n,\nabla\mathcal{L}_n;\eta_Q), &
  \theta_n^-&\leftarrow(1-\tau)\theta_n^-+\tau\theta_n.
  \label{eq:polyak}
\end{align}
The implemented Huber threshold is one in normalized network-output units;
gradients are clipped before the Adam update. All structurally valid DQN replay
samples may update the Q network; the old PPO contention gate, probability-ratio
clipping, and actor--critic advantage objective do not apply.

\paragraph{Offline-to-online replay.}
Exploration data provide initialization and continue to be mixed with online
experience. For $M_n^{\mathrm{on}}$ online decisions and an offline pool of size
$M_n^{\mathrm{pool}}$, the number sampled from the offline pool is
\begin{equation}
  M_n^{\mathrm{off}}=\min\{M_n^{\mathrm{pool}},
              \max(\lfloor\zeta M_n^{\mathrm{on}}\rfloor,M_{\min}^{\mathrm{off}})\}.
  \label{eq:offline-mix}
\end{equation}
Thus $\zeta$ is an offline-to-online count ratio, not directly the offline
fraction of the combined batch. Action holds and returns are prepared
consistently for both sources. Invalid actions or incompatible state encodings
are excluded. The same preprocessing is used by the Local training service and
FL clients. Learning is off-policy, but adequate coverage of meaningful actions
and demand/channel conditions is still needed to limit unsupported Q estimates.

\section{Sample-Weighted Federated Averaging}
\label{subsec:fl}

\paragraph{Participants and local computation.}
Federated rounds are launched every $\Delta_{\mathrm{FL}}$ simulated seconds.
The federated clients are the IAB-node agents $\Nset$, excluding the donor.
All use the same Q-network architecture, feature meanings, normalization, mask
menu, reward, discount, and replay processing. Let
$\Theta_n=(\theta_n,\theta_n^-)$ include the online and target parameters; the
legacy checkpoint names \texttt{actor} and \texttt{critic} refer to these two
Q networks in DQN mode, not to an actor--critic algorithm.
At round $r$, participating client $n\in\mathcal{P}_r\subseteq\Nset$ receives
$\Theta_{\mathrm{glob}}^{(r-1)}$ and performs $E_{\mathrm{FL}}$ local DQN updates
on its processed dataset $\mathcal{D}_n^{(r)}$:
\begin{equation}
  \Theta_n^{(r)}=\operatorname{LocalDQN}_{E_{\mathrm{FL}}}
                 (\Theta_{\mathrm{glob}}^{(r-1)},\mathcal{D}_n^{(r)}).
\end{equation}
Background Local training also continues between federated rounds. Local
optimizers and training counters are retained; they are not averaged.

\paragraph{Aggregation and broadcast.}
Let $w_n^{(r)}$ be the number of valid training decision samples in the client's
training partition after preprocessing, including any offline samples actually
included there. It is neither an equal-node weight nor a count restricted to
PPO-contended samples, and it is not multiplied by the number of update epochs.
A client that cannot perform a valid training update reports zero weight.
For $W^{(r)}=\sum_{n\in\mathcal{P}_r}w_n^{(r)}>0$, standard FedAvg is
\begin{equation}
  \Theta_{\mathrm{glob}}^{(r)}=
        \sum_{n\in\mathcal{P}_r}\frac{w_n^{(r)}}{W^{(r)}}\Theta_n^{(r)}.
  \label{eq:fedavg}
\end{equation}
Both online and target networks are averaged with these weights. The result is
published to all configured clients, including zero-weight clients, and becomes
available to inference through checkpoint reload. If $W^{(r)}=0$, aggregation is
skipped and no new averaged model is broadcast. Clients retain their current
local state. Aggregation and reload introduce delay; nodes need not hold
identical parameters between rounds because they continue local updates.

FedAvg exchanges parameters and sample counts rather than raw local replay
records through the FL protocol. The experimental services are nevertheless
co-located on PC1 and use shared database infrastructure. This deployment does
not establish physical data isolation, cryptographic privacy, or an additional
privacy guarantee. FedAvg does not directly distribute radio slots or PRBs,
and its potential goodput benefit is an empirical question under heterogeneous
node data.

\paragraph{Controlled comparison.}
Stage 1 is PF, Stage 1.5 is pure Local DQN, and Stage 2 adds FedAvg to that same
Local design. The Local architecture, initialization, offline data, optimizer
initial state, observation features, actions, reward, exploration schedule,
background-update settings, and training-seed schedule are held identical.
Stage 1.5 restarts from the common initial state and does not inherit Stage 2's
online replay or trained models. FedAvg contributes additional client updates
as well as parameter sharing; both the Local and FL update counts must be
reported, and equal-computation controls are needed to isolate sharing from
extra optimization work. Checkpoint selection uses validation data, with final
performance measured on seeds excluded from tuning and with repeated runs.

\section{Mapping to the O-RAN Architecture}
\label{subsec:oran}

The platform contains OAI E2 nodes and a FlexRIC Near-RT RIC. Local xApps collect
MAC statistics, relay observations through ZeroMQ, and enforce the returned
per-child masks using the custom MAC control path. A Python service called
``Local rApp'' supplies DQN inference and background training. The Flower server
and clients form the service called ``Global rApp''. These names denote custom
services in this implementation; no Non-RT RIC, A1 interface, or SMO is deployed,
so these services are not claimed to be standard Non-RT RIC rApps. Global xApp
observation broadcasting and Global rApp federated aggregation are distinct
functions.

\paragraph{Time conversion.}
With simulator speed $S_{\mathrm{sim}}$, elapsed time and measured rates obey
\begin{equation}
  T_{\mathrm{sim}}=S_{\mathrm{sim}}T_{\mathrm{wall}},\qquad
  R_{\mathrm{sim}}=R_{\mathrm{wall}}/S_{\mathrm{sim}}.
  \label{eq:sim-time}
\end{equation}
Software timers are configured in wall-clock units and converted by this rule;
radio slot duration is already a simulated-time quantity and is not multiplied
again. The current setting is $S_{\mathrm{sim}}=0.3$.

\begin{table}[htbp]
\centering
\small
\caption{Current deployment periods, expressed in simulated time.}
\label{tab:oran-mapping}
\begin{tabular}{@{}p{0.31\textwidth}p{0.40\textwidth}p{0.19\textwidth}@{}}
\toprule
Function & Implementation & Period \\
\midrule
Slot-level PF & OAI DU MAC & 0.5 ms \\
MAC indication & Custom E2SM-MAC report & 0.03 s \\
Observation and control refresh & Local xApp / ZeroMQ & 0.3 s \\
Joint-action decision & Local DQN service & 1.5 s \\
Background Local training & Local DQN service & 18 s \\
Global observation broadcast & Global xApp, when enabled & 0.6 s \\
FedAvg round launch & Flower scheduler, Stage 2 & 54 s \\
Checkpoint reload polling & Local DQN service & 9 s \\
\bottomrule
\end{tabular}
\end{table}
These are configured periods, not upper bounds on end-to-end completion latency.
The ZMQ timeout is a wall-clock software deadline; it is not a radio-slot deadline.
If communication fails, the xApp requests PF fallback. A stalled control process
can leave the previous mask stored in the MAC, so actual enforcement and phase
transitions must be checked independently of inference decisions.

\paragraph{Current numerical instance.}
The deployed numerology is $\mu=1$, with $B_n=106$ PRBs and a 10-slot TDD period:
seven full downlink slots, one special slot with six downlink symbols, and two
uplink slots. The mask width is $W=16$. The MAC uses
$\iota_n(t)=s_n(t)\bmod16$, where $s_n(t)$ is the slot index within a radio frame;
it does not use the absolute slot index modulo 16. The finite menu is
\begin{equation}
  \mathcal{M}=\{\mathtt{0x0101},\mathtt{0x1111},\mathtt{0x9044},
                 \mathtt{0x9292},\mathtt{0xab98},\mathtt{0xffff}\}.
\end{equation}
The current formulation uses $L_h=5$, $\gamma=0.5$, $n_{\mathrm{step}}=3$,
$\alpha=0$, $\tau=0.01$, $c_Q=50$, and minibatches of 128 decisions. The flat
Q network has two 128-unit hidden layers. Learning rate, exploration parameters,
PF margin, offline-mixing ratio, and Local/FL update counts are run-specific
settings and must be reported from the executed configuration and checkpoints,
not inferred from source-code defaults. Learning is disabled and exploration is
zero during final evaluation.

\paragraph{Scope of generalization.}
Topology size enters the equations through sets and descendant relationships,
not through fixed node numbers. Nevertheless, the present code uses a padded
capacity $C_{\max}=16$ children per DU and explicit mappings for the experimental
parent/child relationships and reward joins. A larger or deeper deployment
requires configuring those relationships, recursively collecting descendant
terminal measurements, and retaining correct MT/UE identities. If a DU exceeds
$C_{\max}$, the current flat encoding and network dimensions must be expanded
and compatible models prepared; the present implementation must not be treated
as supporting that case unchanged.

For any finite deployment one can choose
$C_{\max}\ge\max_n|\Cset_n|$ and retain a common input layout for FedAvg.
However, enumerating~\eqref{eq:action-count} costs exponentially many candidates
in the number of attached direct UEs, so increasing UE count can require
candidate pruning or another structured action search. Generalized notation does
not prove scalability, permutation robustness, or transferable goodput gains.
Different depths, fan-outs, demand levels, channel conditions, and feedback
configurations require separate evaluation.

\section{Notation}
\label{subsec:notation}

\begin{longtable}{@{}p{0.28\textwidth}p{0.66\textwidth}@{}}
\toprule
Symbol & Meaning \\
\midrule
\endhead
$\Nset,\Uset,\parent(v)$ & IAB-node set, terminal-UE set, serving parent \\
$\Cset_n^{\mathrm{I}},\Cset_n^{\mathrm{U}}$ & IAB children (MT entities), directly attached UEs \\
$h(v),H,\mathcal{P}(u),\Uset_n$ & Depth, maximum UE hop count, downlink path, subtree terminal UEs \\
$\Delta_o,\Delta_d,L_h,\Tset_k$ & Observation interval, normal decision interval, hold length, decision slots \\
$B_n,\tilde B_n,\beta_n$ & Total PRBs, available PRBs, available fraction \\
$b_{n,v},z_{n,v},e_{n,v}$ & Total allocated PRBs, new-transmission indicator, MAC eligibility \\
$\rho_{n,v},\eta_{n,v}$ & Payload capacity per PRB, effective delivery fraction \\
$\mathcal{M},W,\iota_n(t)$ & Mask menu, mask width, MAC mask-bit index \\
$\Aset_n,a_{\mathrm{PF}},\mathcal{S}_n$ & Valid joint actions, all-open action, selected direct UEs \\
$Q_{n,v},A_{n,v},S_{n,v},D_{n,v}$ & Aggregate queue, arrival bits, forwarded bits, dropped bits \\
$\lambda_u,x_u^{\mathrm{app}},\widehat x_u$ & Offered rate, application goodput, MAC throughput proxy \\
$\bm z_{n,v},\bm g_n,\hat p_n,\hat c_n$ & Child features, node context, upstream/downstream congestion proxies \\
$C_{\max},Y_{\mathrm{ref}},Q_{\mathrm{ref}}$ & Padded child capacity, byte-count normalization references \\
$r_n,\overline r_n,G_{n,k}^{(m)}$ & Observation reward, hold-averaged reward, multi-step return \\
$\gamma,n_{\mathrm{step}},\delta,c_Q$ & Per-decision discount, return length, PF margin, Q output scale \\
$\theta_n,\theta_n^-,\tau,\eta_Q$ & Online/target parameters, Polyak factor, learning rate \\
$\zeta,M_n^{\mathrm{off}}$ & Offline-to-online count ratio, sampled offline decisions \\
$\Theta_n,w_n^{(r)},\mathcal{P}_r$ & FL parameter pair, valid training-sample weight, responding clients \\
$\Delta_{\mathrm{FL}},S_{\mathrm{sim}}$ & Federated launch interval, simulator speed factor \\
\bottomrule
\end{longtable}

\end{EnChapter}
